% Paper 5, §13 — Scalable alignment: the oversight layer
% Anchors: kb/notes/alignment/oversight-and-scalable-alignment.md
%          kb/excerpts/irving2018-debate.md
%          kb/excerpts/burns2023-w2s.md (cross-paper, Paper 2 §weak-to-strong)

\section{Scalable alignment: the oversight layer}
\label{sec:scalable-oversight}

The threats covered in
\cref{sec:scheming,sec:alignment-faking,sec:watermarking} are
deployment-distribution problems on \emph{current} models: \glsterm{scheming}{scheming}
under elicited goals, alignment-faking under training distribution
shift, \glsterm{watermark}{watermark} survival under removal pressure. The threat in this
section is one layer further out --- the question of whether any
training or evaluation protocol the field has yet proposed can
\emph{align} a model whose capability already exceeds the human
overseer's, using a bounded human resource budget. The thesis of this
section is that \citet{irving2018-debate} supplies the canonical
theoretical argument that such a protocol can in principle exist ---
\glsterm{debate}{debate} between two AI agents, judged by a weaker human, decides
$\mathbf{PSPACE}$ under polynomial-time judging --- but that the 2024
\glsterm{llm}{LLM}-\glsterm{debate}{debate} empirical literature gives \emph{partial} support in
verifiable-claim domains, \emph{no clear} support in open-ended ones,
and that as of 2026-Q2 no candidate protocol --- \glsterm{debate}{debate}, iterated
distillation/amplification, recursive reward modeling, weak-to-strong
generalization, or constitutional approaches --- is known to scale to
\emph{superhuman} capability with a \emph{bounded} human budget. The
gap between the $\mathbf{PSPACE}$ result and the deployment-grade
protocol is the alignment-relevant version of the
$\mathbf{P}$-vs-$\mathbf{PSPACE}$ question, and it is open.

\subsection{The hierarchy of alignment difficulty}
\label{sec:oversight-hierarchy}

The motivating framing for the entire section is Irving et al.'s
hierarchy of alignment-difficulty levels, defined by what the human
overseer can do with a given task%
\footnote{KB excerpt: \texttt{kb/excerpts/irving2018-\glsterm{debate}{debate}.md\#sec-1}.}
\citep[\S 1]{irving2018-debate}:

\begin{quote}\itshape
For some tasks it is harder to bring behavior in line with human goals
than for others. In simple cases, humans can directly demonstrate the
behavior --- this is the case of supervised learning or imitation
learning. \ldots\ Taking a step up in alignment difficulty, some tasks
are too difficult for a human to perform, but a human can still judge
the quality of behavior or answers once shown to them \ldots\ This is
the case of human preference-based reinforcement learning. We can make
an analogy between these two levels and the complexity classes
$\mathbf{P}$ and $\mathbf{NP}$: answers that can be computed easily
and answers that can be checked easily.
\end{quote}

The third level is the regime this section is about
\citep[\S 1]{irving2018-debate}:

\begin{quote}\itshape
A human may be unable to judge whether an explained answer or
exhibited behavior is correct: the behavior may be too hard to
understand without help, or the answer to a question may have a flaw
that is too subtle for the human to detect. We could imagine a system
trained to both give answers and point out flaws in answers; this gives
a third level of difficulty. Flaws themselves may be too hard to judge:
flaws could have their own flaws that must be pointed out to a human.
And flaws of flaws can have flaws, etc. \ldots\ This hierarchy of
alignment tasks has a natural limit: a \glsterm{debate}{debate} between competing agents
where agents make arguments, other agents poke holes in those arguments,
and so on until we have enough information to decide the truth.
\end{quote}

Three observations frame the rest of the section. First, the hierarchy
is structured by judging cost, not by capability axis: each level
relaxes a different constraint on the human's ability to verify a
claim, not on the model's ability to make one. Second, the recursion
``flaws of flaws of flaws'' is the load-bearing protocol primitive: a
\glsterm{debate}{debate}'s job is to drive the dispute to a level the judge \emph{can}
adjudicate, not to render the entire claim transparent. Third, the
hierarchy connects directly to complexity theory: $\mathbf{P}$ is the
demonstration regime, $\mathbf{NP}$ is the preference-learning regime
(judge a single witness), and the next level up is what \glsterm{debate}{debate} is
designed to reach.

\subsection{The debate game}
\label{sec:oversight-debate-game}

The \glsterm{debate}{debate}-game protocol is defined verbatim%
\footnote{KB excerpt: \texttt{kb/excerpts/irving2018-\glsterm{debate}{debate}.md\#sec-2}.}
\citep[\S 2]{irving2018-debate}:

\begin{quote}\itshape
We have a set of questions $Q$, answers $A$, and \glsterm{debate}{debate} statements $S$.
The simplest version of \glsterm{debate}{debate} has two agents competing to convince a
human judge:
\begin{enumerate}
\item A question $q \in Q$ is shown to both agents.
\item The two agents state their answers $a_0, a_1 \in A$ (which may
be the same).
\item The two agents take turns making statements
$s_0, s_1, \ldots, s_{n-1} \in S$.
\item The judge sees the \glsterm{debate}{debate} $(q, a, s)$ and decides which agent
wins.
\item The game is zero sum: each agent maximizes their probability of
winning.
\end{enumerate}
\end{quote}

The thesis the protocol is designed to operationalize is the central
claim%
\footnote{KB excerpt: \texttt{kb/excerpts/irving2018-\glsterm{debate}{debate}.md\#sec-2}.}
\citep[\S 2, central claim]{irving2018-debate}:

\begin{quote}\itshape
\textbf{Claim.} In the \glsterm{debate}{debate} game, it is harder to lie than to refute
a lie. \ldots\ Whether this claim is true for any particular setting
is empirical, though we give some evidence for it below.
\end{quote}

The claim is what does the alignment work. If lying is structurally
more expensive than refuting a lie, then the truthful debater wins more
often than the deceptive one under optimal play, and the judge does
not need to verify the full argument --- only to recognize the better
defense. The protocol's design choice is the bounded-length game: at
each turn a debater pays the cost of one statement to either narrow
the dispute by one level or to commit to a falsifiable claim that the
opponent can challenge%
\footnote{KB excerpt: \texttt{kb/excerpts/irving2018-\glsterm{debate}{debate}.md\#sec-2-1}.}
\citep[\S 2.1, prime-counting recursion]{irving2018-debate}: a
debater lying about $\pi(x)$ is forced to lie about one of $\pi(x/2)$
or $\pi(x) - \pi(x/2)$, halving the dispute's search space per round.

\subsection{The PSPACE construction}
\label{sec:oversight-pspace}

The complexity-class progression that the abstract advertises is
constructed step-by-step%
\footnote{KB excerpt: \texttt{kb/excerpts/irving2018-\glsterm{debate}{debate}.md\#sec-2-2}.}
\citep[\S 2.2, p.~5]{irving2018-debate}:

\begin{quote}\itshape
If we use $H$ to answer questions directly as $H(q)$, we get the
complexity class $\mathbf{P}$ of polynomial time algorithms. \ldots\
If we use the polynomial time algorithm not to output answers but to
judge them, we get the complexity class $\mathbf{NP}$ of questions
with polynomial time checkable witnesses. Instead of $H(q)$ we output
$\exists x. H(q, x)$ where $x$ is a witness found by the ML model.
$\mathbf{NP}$ corresponds to single agent reinforcement learning: an
agent with sufficient capacity can solve tasks the human cannot, but
the human must be able to judge whether the solution is good.

Now consider an adversarial \glsterm{debate}{debate} of length two, where Alice chooses
$x$ attempting to make the human say yes, Bob chooses $y$ attempting
to make the human say no, and the human decides who is correct. The
result is $\exists x \forall y. H(q, x, y)$. \ldots\ This complexity
class is $\Sigma_2 \mathbf{P}$, two steps up the polynomial hierarchy.

We can continue this process for any number of rounds, with Alice and
Bob alternating points and counterpoints, producing the formula
$\exists x_0 \forall x_1 \cdots \exists x_{n-1}. H(q, x_0, \ldots)$
for $n$ rounds of \glsterm{debate}{debate}. \ldots\ If the number of rounds $n$ is
allowed to grow polynomially in the size of the question $q$, the
complexity class is $\mathbf{PSPACE}$: all questions decidable by
polynomial space algorithms.
\end{quote}

The headline theorem%
\footnote{KB excerpt: \texttt{kb/excerpts/irving2018-\glsterm{debate}{debate}.md\#sec-2-2}.}
\citep[Theorem 1, \S 2.2]{irving2018-debate}:

\begin{quote}\itshape
\textbf{Theorem 1.} For any problem $L \in \mathbf{PSPACE}$, we can
choose a polynomial time judge so that optimal play in the \glsterm{debate}{debate} game
with polynomially many steps solves $L$.
\end{quote}

Three structural facts make the construction load-bearing. First, the
debaters are computationally unbounded relative to the judge --- the
construction quantifies over ``optimal play,'' which lets the debaters
play strategies of arbitrary cost in the game tree, while the judge
runs in polynomial time. Second, the alternation $\exists \forall
\cdots$ is what walks up the polynomial hierarchy, and growing the
number of rounds polynomially is what closes it into $\mathbf{PSPACE}$.
Third, the equilibrium-strategy lemma underlying the theorem is that,
under optimal play of a zero-sum game with a fixed judge, the truthful
debater wins iff the truthful answer is the one the judge would output
given access to the full game tree --- which is exactly the
$\mathbf{PSPACE}$ \glsterm{value}{value} of the underlying alternating-quantifier
formula.

The interpretive caveat is in the same paragraph%
\footnote{KB excerpt: \texttt{kb/excerpts/irving2018-\glsterm{debate}{debate}.md\#sec-2-2}.}
\citep[\S 2.2]{irving2018-debate}:

\begin{quote}\itshape
These complexity class arguments are analogies only: we do not expect
tractable machine learning algorithms to achieve all of $\mathbf{PSPACE}$.
Rather, the analogies show that at least in theory we can be limited
only by the capacity of the ML models and our ability to train them,
not the supervisory signal. This gives us hope that \glsterm{debate}{debate} could
resolve AI alignment without sacrificing model strength.
\end{quote}

\begin{analogy}
The \glsterm{debate}{debate}-game theorem is the alignment-protocol cousin of an
\emph{interactive proof system}. In the IP framework, an unbounded
prover convinces a polynomial-time verifier of membership in a
language $L$ by exchanging messages, and the celebrated result
$\mathbf{IP} = \mathbf{PSPACE}$ \deepencite{shamir1992-ip,
citation-pending} says that interaction lifts the verifier from
$\mathbf{NP}$ to $\mathbf{PSPACE}$ exactly when the prover is allowed
arbitrary computation. \glsterm{debate}{Debate} adds one structural ingredient ---
\emph{adversarial} provers in a zero-sum game --- which is the dual
construction: where IP's single prover plays cooperatively to convince
the verifier of a true claim, \glsterm{debate}{debate}'s two provers play competitively
and the equilibrium-honesty lemma claims that truth is the strategy
that wins. (Returning to canonical form: the alternating-quantifier
formula $\exists x_0 \forall x_1 \cdots H(q, x_0, \ldots)$ in
\cref{sec:oversight-pspace} is the same game-tree object that the IP
construction unrolls; \glsterm{debate}{debate} fixes the strategy of one prover at
``honest'' and asks whether the resulting equilibrium recovers the
formula's true \glsterm{value}{value}.)
\end{analogy}

The MNIST sparse-classifier experiment in the original paper is the
toy empirical confirmation%
\footnote{KB excerpt: \texttt{kb/excerpts/irving2018-\glsterm{debate}{debate}.md\#sec-3-1}.}:
two debaters reveal pixels in turn to a sparse classifier judge,
and the truthful debater's pixel choice boosts judge accuracy from
$59.4\%$ to $88.9\%$ at $6$ pixels and from $48.2\%$ to $85.2\%$ at
$4$ pixels \citep[\S{}abstract,\S 3.1]{irving2018-debate}. The judge
sees the same number of pixels in both conditions; the \glsterm{debate}{debate}
condition wins because the truthful debater chooses informative
pixels.

\subsection{2024 LLM-debate empirics: partial support}
\label{sec:oversight-llm-debate}

The 2024 \glsterm{llm}{LLM}-\glsterm{debate}{debate} literature is the first round of empirical work
that operationalizes Irving et al.'s claim at language-model scale.
\citet{kenton2024-debate-llm-judges} (Kenton et al., NeurIPS 2024)
generalize the Irving 2018 setup to \glsterm{llm}{LLM} debaters and an
\emph{intentionally weaker} \glsterm{llm}{LLM} judge, comparing three conditions ---
\glsterm{debate}{debate} (two debaters, opposing positions), \glsterm{consultancy}{consultancy} (single
debater), and direct QA (judge alone)%
\footnote{KB note: \texttt{kb/notes/alignment/oversight-and-scalable-alignment.md\#2-4-\glsterm{debate}{debate}-with-weak-\glsterm{situational-awareness}{llm}-judges-kenton-et-al-2024}.}
\deepencite{kenton2024-debate-llm-judges §abstract,
citation-pending}. The headline findings are mixed: \glsterm{debate}{debate}
\emph{outperforms} \glsterm{consultancy}{consultancy} across all task domains tested; \glsterm{debate}{debate}
versus direct QA is task-dependent --- \glsterm{debate}{debate} wins on
information-asymmetry QA where the debaters' access to a long passage
gives them a verifiable claim to defend, and produces mixed results on
math, coding, logic, and multimodal where the claims being defended
are open-ended. \citet{khan2024-debating-llms} (Khan et al., 2024)
report a complementary result --- ``debating with more persuasive LLMs
leads to more truthful answers'' \deepencite{khan2024-debating-llms
§main-result, citation-pending} --- in a setup where the ground truth
is fixed externally, so the persuasiveness gradient lines up with the
truth gradient.

The Brown-Cohen et al.\ 2024 \emph{doubly-efficient \glsterm{debate}{debate}} result
\deepencite{brown-cohen2024-doubly-efficient-debate §main-theorem,
citation-pending} is the strongest current bridge between the
$\mathbf{PSPACE}$ argument and an empirically testable protocol:
optimizing for \emph{persuasiveness} in \glsterm{debate}{debate} \emph{does} drive
truth-finding under specific conditions, reported at $76$--$88\%$
accuracy in the cited regime (pending direct excerpt verification)%
\footnote{KB note: \texttt{kb/notes/alignment/oversight-and-scalable-alignment.md\#2-5-doubly-efficient-\glsterm{debate}{debate}-brown-cohen-et-al-2024}.}.
The mechanism --- pending direct excerpt --- is that an honest
prover's strategy in the doubly-efficient construction is structurally
cheaper than a dishonest prover's, so optimization pressure on
persuasiveness recovers the equilibrium-honesty claim that Irving
2018 left empirical.

The split that the Kenton 2024 results expose is the diagnostic one
for the section. In domains where the ground truth is a verifiable
claim about a specified artifact --- a passage, a fixed reference, an
externally-graded answer --- the \glsterm{debate}{debate} condition lifts judge
accuracy above direct QA, in line with Irving's prediction. In domains
where the ground truth is open-ended --- math at the open-research
boundary, code with no test suite, philosophical or aesthetic claims
--- the same condition does not yield a consistent gain. The 2024
literature is partial support for the Irving prediction, restricted
to the regime where the dispute can be driven down to a verifiable
node in finite rounds.

\subsection{Adjacent protocols: IDA, RRM, W2S, constitutional}
\label{sec:oversight-adjacent}

\glsterm{debate}{Debate} is one candidate scalable-oversight protocol among several,
each with a different decomposition target%
\footnote{KB note: \texttt{kb/notes/alignment/oversight-and-scalable-alignment.md\#6}.}.
\textbf{Iterated distillation/amplification} (\glsterm{iterated-amplification}{IDA}, Christiano 2018)
\deepencite{christiano2018-ida §abstract, citation-pending} runs the
recursion $A_{k+1} = \mathrm{Distill}(\mathrm{Amplify}(A_k))$ where
Amplify decomposes a hard task into subtasks small enough for $A_k$
to solve and Distill trains a single network to imitate Amplify's
composed answer; the supervision signal at each step is grounded in
the human-level $A_0$ via the recursion%
\footnote{KB note: \texttt{kb/notes/alignment/oversight-and-scalable-alignment.md\#3-iterated-amplification--\glsterm{iterated-amplification}{ida}-christiano-2018}.}.
No scaled empirical instantiation in modern LLMs has been published
as a clean test; the conceptual machinery has been partially absorbed
into \glsterm{rlaif-3}{RLAIF} and into \glsterm{debate}{debate} itself. \textbf{Recursive reward modeling}
(Leike et al.\ 2018) \deepencite{leike2018-rrm §abstract,
citation-pending} is the dual: rather than recursing on the task, it
recurses on the \emph{reward function}, training reward models that
themselves use AI-assistant outputs to evaluate hard cases.

\textbf{\glsterm{constitutional-ai}{Constitutional AI}} \citep{bai2022-cai} decomposes the
\emph{\glsterm{value}{value} spec} into a \glsterm{constitution}{constitution} of natural-language principles,
then uses the model itself to critique-and-revise outputs against the
\glsterm{constitution}{constitution}; the supervision signal is concentrated in the
\glsterm{constitution}{constitution}, downstream labels are AI-generated. \glsterm{constitutional-ai}{CAI} is the only
candidate in this list with production deployment (Anthropic \glsterm{claude}{Claude}
family), but its theoretical leverage decays as model capability
exceeds the \glsterm{constitution}{constitution}-author's ability to write a comprehensive
spec --- the \glsterm{constitution}{constitution} caps the supervision precision%
\footnote{KB note: \texttt{kb/notes/alignment/oversight-and-scalable-alignment.md\#4-3-where-it-fits-in-the-scalable-oversight-space}.}.

\textbf{Weak-to-strong generalization}
\citep{burns2023-w2s}\footnote{Cross-paper: see also Paper~2
\S{}weak-to-strong (Burns 2023).} decomposes the supervision
\emph{signal} itself: train a strong student on labels generated by a
weak supervisor, measure the \glsterm{performance-gap-recovered}{Performance Gap Recovered} (\glsterm{performance-gap-recovered}{PGR})
\[
\mathrm{PGR} \;=\; \frac{c_{\text{strong-student}} - c_{\text{weak-supervisor}}}{c_{\text{strong-ceiling}} - c_{\text{weak-supervisor}}}
\]
where $c_{\text{strong-ceiling}}$ is the strong model trained on
ground-truth labels%
\footnote{KB note: \texttt{kb/notes/alignment/oversight-and-scalable-alignment.md\#5-weak-to-strong-generalization-burns-et-al-2023}.}.
With an \glsterm{auxiliary-confidence-loss}{auxiliary confidence loss}, the GPT-2 $\to$ GPT-4-family
experiments recover ``close to GPT-3.5-level'' on NLP tasks, less on
chess and reward modeling. W2S is methodologically the cleanest
deployable instrument in this list --- you can run it today on
checkpoints you have --- but it is also the construction whose
analogy-to-superhuman-supervision is the most contested: the
GPT-2/GPT-4 capability gap may be structurally dis-analogous to a
human-vs-superhuman gap, since both ends share training distribution
and the same regime of representable concepts%
\footnote{KB note: \texttt{kb/notes/alignment/oversight-and-scalable-alignment.md\#5-3-why-this-is-the-methodology-not-the-answer}.}.

\textbf{Process supervision and PRMs}, covered in Paper~3
\S{}process-supervision-and-PRMs, supplies an adjacent training-side
signal: rather than rewarding final answers, reward each
\emph{intermediate} reasoning step. The \glsterm{process-reward-model}{PRM} is itself a learned
judge, trained on human or AI process annotations, and its operating
point is the same as the \glsterm{debate}{debate}-judge's: a polynomial-cost grader
selecting among candidate continuations from a stronger reasoner.
Where \glsterm{debate}{debate} is an inference-time protocol and W2S is a training-time
one, PRMs straddle: trained once on process labels, deployed
inference-time as a step-by-step reward. The four candidate protocols
are non-mutually-exclusive --- the speculative composition flagged in
the KB note%
\footnote{KB note: \texttt{kb/notes/alignment/oversight-and-scalable-alignment.md\#7}
\S 7 (\textsc{speculation}).} is \glsterm{debate}{debate} or W2S as the
\emph{selection} mechanism with interpretability as the
\emph{verification} layer; no published end-to-end pipeline
implements this composition as of 2026-Q2.

\begin{contradiction}
\textbf{\glsterm{debate}{Debate}-as-alignment-protocol viability.} The Irving 2018
$\mathbf{PSPACE}$ result \citep[Theorem 1]{irving2018-debate} is a
theoretical existence claim under an equilibrium-strategy assumption:
optimal play recovers the alternating-quantifier \glsterm{value}{value}, and the
truthful strategy wins. The 2024 \glsterm{llm}{LLM}-\glsterm{debate}{debate} empirical work is
\emph{partial} support: \citet{kenton2024-debate-llm-judges} report
\glsterm{debate}{debate} beats direct QA on information-asymmetry QA where the dispute
is over a verifiable claim, but is mixed-or-flat on math, code, and
open-ended domains; the Brown-Cohen 2024 doubly-efficient construction
\deepencite{brown-cohen2024-doubly-efficient-debate §main-theorem,
citation-pending} hits $76$--$88\%$ accuracy in the reported regime
but the regime conditions are restrictive. Whether equilibrium
honesty holds in practice --- whether self-play under bounded compute
on real LLMs converges to the Irving 2018 honest equilibrium rather
than to a dishonest one that exploits the judge's weakness --- is
open. The sharpest unresolved test is the one this section's
contradiction in \cref{sec:contradictions} points to: an empirical
demonstration of \glsterm{debate}{debate}-induced honesty on a domain where the judge
is \emph{strictly weaker} than both debaters and where the question
is \emph{not pre-verifiable}. The current empirical work does not
meet both conditions.
\end{contradiction}

\subsection{The bounded-budget question}
\label{sec:oversight-bounded-budget}

The constraint that organizes all four protocols is the
\emph{bounded human budget}: the field's claim is not that
human-judging cost is irrelevant but that it should grow
\emph{sub-linearly} in the capability being supervised. None of \glsterm{debate}{debate},
\glsterm{iterated-amplification}{IDA}, RRM, W2S, or constitutional approaches has yet been demonstrated
to scale to \emph{superhuman} capability with a bounded human resource
budget; this is the empirical version of
$\mathbf{P}$-vs-$\mathbf{PSPACE}$, with two structural caveats. First,
the complexity-theoretic argument is asymptotic and unconditional on
the equilibrium-strategy assumption; the empirical version requires
the equilibrium to be reachable by self-play on bounded-capacity
debaters. Second, the alignment-relevant variant of the question
includes adversarial robustness against models that can recognize the
oversight protocol --- a form of strategic awareness that the
$\mathbf{PSPACE}$ construction does not contemplate, but that the
\glsterm{scheming}{scheming} results of \cref{sec:scheming} demonstrate is non-zero in
2024-frontier models.

\begin{intuition}
The relationship between \glsterm{debate}{debate}'s $\mathbf{PSPACE}$ theorem and a
deployment-grade alignment protocol is the same shape as the
relationship between a worst-case complexity result and a practical
algorithm. The theorem says \emph{the supervisory signal is not the
bottleneck}: a polynomial-time judge plus optimally-playing debaters
can in principle decide any $\mathbf{PSPACE}$ language, so the
bottleneck must lie in the trainability of the debaters or the
realizability of optimal play under bounded compute. The empirical
program of the last three years has been to translate that ``in
principle'' into ``in practice on language models,'' and the result
so far is that the in-principle ceiling is far above what current
debater-training procedures and judge architectures actually reach.
The same shape recurs in interactive-proof theory: $\mathbf{IP} =
\mathbf{PSPACE}$ is unconditional, but the protocols that achieve it
in deployment are restricted to specific languages with specific
arithmetization tricks, not to arbitrary $\mathbf{PSPACE}$ languages.
(Returning to canonical form: the gap between the alternating-quantifier
formula in \cref{sec:oversight-pspace} and a specific
\glsterm{debate}{debate}-trained-on-LLMs implementation is the gap between
``optimal-play'' and ``self-play with finite compute and finite
debater capacity''; the empirical literature is the search for the
class of formulae for which the gap closes.)
\end{intuition}

The shared structural limit across all four protocols is that they
are \emph{outer-loop} procedures: they observe the model's behavior,
not its internal computation. A model capable of strategic deception
that recognizes the oversight protocol can in principle defeat each
of them by behaving aligned to the weak supervisor / \glsterm{debate}{debate} judge /
\glsterm{constitution}{constitution}-grader. This is the
\cref{sec:scheming,sec:alignment-faking} threat re-encountered at the
oversight layer: the same elicited capability that
\citet{meinke2024-apollo-scheming} demonstrate in agentic scaffolds
is exactly what a sufficiently situationally-aware debater can deploy
against a weaker judge. The natural complement is interpretability
access, which Paper~4 \S{}probing-and-circuits treats as an
internal-verification channel: probes for ``model believes it is in
an oversight protocol'' or activation-patching for goal-tracking
features. As of 2026-Q2 no published end-to-end pipeline composes a
behavioral oversight protocol with an interpretability-grounded
verification layer; this is the most-cited next step in the
scalable-oversight literature%
\footnote{KB note: \texttt{kb/notes/alignment/oversight-and-scalable-alignment.md\#7}
\S 7 (\textsc{speculation}).}.

\subsection{Forward link to the contradictions discussion}
\label{sec:oversight-forward}

\glsterm{debate}{Debate} is one protocol among several, and the field as of 2026-Q2 has
none with bounded human cost at superhuman capability. The
contradiction surfaced in this section --- the gap between Irving's
$\mathbf{PSPACE}$ result and the partial-and-domain-restricted 2024
\glsterm{llm}{LLM}-\glsterm{debate}{debate} empirics --- joins the alignment-evaluation contradictions
of \cref{sec:scheming,sec:alignment-faking} and the \glsterm{watermark}{watermark}-robustness
contradiction of \cref{sec:watermarking} in the concentrated discussion
of \cref{sec:contradictions}. The evidence that would close it is the
one this paper commits to naming: an empirical demonstration of
\glsterm{debate}{debate}-induced honesty on a domain where the judge is strictly weaker
than both debaters \emph{and} where the question is not pre-verifiable.
The current empirical work --- Kenton 2024, Khan 2024, Brown-Cohen
2024 --- meets one or the other condition, not both. The structural
location of this section in the paper's layered-defense view is the
outermost layer: where \cref{sec:adversarial-robustness} measures
attack-elicited misbehavior on current models and
\cref{sec:scheming,sec:alignment-faking} measure capability for
deception in elicited and training-shifted regimes, this section asks
the question that comes after all of those have been answered ---
whether any oversight protocol the field has yet proposed will still
work when the model being overseen is the one we cannot directly
evaluate.
